\section{Verification purpose}

This case verifies the coupled response from Wildland--Urban ELMFIRE (WU--E) source wind and building
geometry through the model-2 ellipse and heat-flux routines to the
transient direct-flame-contact (DFC) and radiation rasters.  It is capable of
exposing an incorrect wind-unit conversion, a wrong branch predicate, a
regression-coefficient or ellipse-construction error, a wind-direction sign or
raster-orientation error, failure of the separation clamp, and disagreement
between constant and raster building-parameter input paths.

One initially burning structure is held stationary by a zero spread-adjustment
raster.  Spotting, terrain, moisture variation, and secondary propagation are
disabled. Twenty-three deterministic simulation conditions exercise both source-wind branch
boundaries, Hamada parameters (A) and (D), four right-angle rotations, and
constant/raster equivalence, plus the intended source fuel-reach factors 1 for
structures and 1/3 for vegetation. The reference calculation implements the intended downscaled
WU--E formulation.  The pinned source is the system under test, so a mismatch
is reported rather than copied into the expected result.

\textbf{Verification question.} Do all required simulation conditions reproduce the
intended heat-release rate (HRR) per unit area (HRRPUA) and full-domain DFC/radiation maps within normalized
$L_1$ error \HeatMapTolerance, preserve the paired clamp, rotation, and
constant/raster invariants, and create no additional burning cell?

\ReportOverview

\section{Mathematical formulation and numerical implementation}

The source converts the 20-ft input wind from miles per hour to metres per
second as $V=0.447W$.  The regression is linear when $V<10$, quadratic when
$10\le V\le17.3$, and linear with a distinct coefficient set when
$V>17.3$.  In each branch it produces downwind, upwind, and sidewind reaches
$(K_d,K_u,K_s)$ from $(V,A,D)$, after enforcing

\begin{equation}
  D_{\mathrm{eff}}=\min(D,\SI{50}{\metre}).
\end{equation}

The stored ellipse quantities are

\begin{align}
 B &= \frac{K_d+K_u}{2}, &
 E &= \min\left(\frac{B}{2},B-K_u\right), &
 A_e &= \frac{K_s}{\sqrt{1-(E/B)^2}}.
\end{align}

Here $B$, $A_e$, and $E$ retain the source nomenclature major,
minor, and eccentricity; all three are lengths.  For comparison, the pinned
source currently constructs the following continuous angular boundary, which
is a documented deviation from the intended three-sector model:

\begin{equation}
 r_e(\theta_f)=
 \frac{0.3K_d(B-E)}{A_e^2}\,
 \frac{A_e^2}{B-E\cos\theta_f},
 \qquad \theta_f=\operatorname{atan2}(\Delta y,\Delta x)
                  -\frac{\pi}{180}(270-WD).
\end{equation}

The intended model uses a three-sector reach and a physical-distance
correction.  With $L_c=\texttt{BLDG\_AREA}$ interpreted as a linear
characteristic building dimension in metres and $s_{sd}$ the separation,

\begin{align}
 \beta_{ssd}&=\min\!\left(1,\frac{s_{sd}+L_c}{\Delta x}\right),\\
 K_\theta&=\begin{cases}B+E,&\cos\theta_f>0.7,\\B-E,&\cos\theta_f<-0.7,\\2A_e,&\text{otherwise},\end{cases}
\end{align}

and the target-cell overlap factors are clipped to $[0,1]$:

\begin{align}
 f_d &= \operatorname{clip}\!\left(
       \frac{K_\theta+\Delta x/2-r\beta_{ssd}}{\Delta x},0,1\right),\\
 f_r &= \operatorname{clip}\!\left(
       \frac{HAZ+\Delta x/2-r\beta_{ssd}c}{\Delta x},0,1\right)(1-f_d).
\end{align}

With $A_s=L_c^2$, $C_h=A_s/[\pi(cB)(cA_e)]$, default $c=1$, target
nonburnable fraction $n_b=0$, and absorptivity $\alpha=0.8$, the expected
fields are

\begin{align}
 q''_{\DFC} &= g_f\eta(W)\,0.65(1-n_b)f_d\,\HRRPUA(t)C_h,\\
 q''_{rad} &= g_f\eta(W)\frac{0.35(1-n_b)\alpha f_r\HRRPUA(t)A_s}
                         {4\pi (r\beta_{ssd}c)^2},\\
 \eta(W)&=\frac{3.25}{1+\exp[-0.05(W-75)]}+0.5,\qquad
 g_f=\min\!\left[\left(\frac{\Delta x}{L_c}\right)^2,1\right].
\end{align}

The physical heat-affecting zone is $HAZ=100$ m and its search radius is
$\lceil HAZ/\Delta x\rceil$. Radiation is complementary to DFC through
$(1-f_d)$. The source instead uses a continuous angular ellipse, 0.3 heat
factors, cell-area normalization, and a cell-count band; these are explicit
deviation targets. Preprocessing assigns the prescribed parameter values
independently to each simulation condition.

\section{Derivation of expected behavior}

The ellipse-regression coefficients and branch predicates were independently
transcribed from the specified WU-E spread-rate formulation. The reviewed
configuration interpretation, spread-rate calculation, and level-set formulation
were identical to the specified ELMFIRE source revision; this identity claim
does not extend to all other model components. Purnomo et al. describe the wind-responsive urban flame ellipse
and heat-transfer coupling in the peer-reviewed WU--E publications
\cite{PurnomoPCI2024,PurnomoIJWF2024}.  The detailed equations above document
the intended implemented model directly. The executable snapshot is evaluated against, rather than substituted
for, this formulation.

At $V=10$ and $V=17.3~\si{\metre\per\second}$, equality belongs to the quadratic
branch.  The immediate neighboring simulation conditions on each side must instead use the
linear low- or high-wind coefficients.  The $D=75~\si{\metre}$ field must be
identical to $D=50~\si{\metre}$.  A raster parameter field equal to the constants
must produce the same heat maps.  With all other inputs invariant, rotating
the wind-from direction by $90^\circ$ must rotate the complete DFC-plus-
radiation field by the same angle on this square grid.

The HRRPUA curve rises from zero to $100~\si{\kilo\watt\per\square\metre}$ over one second,
remains at the peak through $t=100~\si{\second}$, decays linearly to zero at
$t=110~\si{\second}$, and stays zero.  The $t=60~\si{\second}$ dump consequently gives
an exact nonzero full-map comparison without conflating ellipse response with
the HRR transient.

\begin{figure}[p]
  \centering
  \includegraphics[width=\textwidth]{expected_response.pdf}
  \caption{Intended-model reference calculation for the branch, characteristic building
  dimension, and separation sweeps. These are predicted values, not ELMFIRE output.}
\end{figure}

\section{Verification metrics and acceptance criteria}

For any reference raster $R$ and ELMFIRE raster $E$, the primary metric is

\begin{equation}
 \nLone=\frac{\sum_{i,j}|E_{ij}-R_{ij}|}{\sum_{i,j}|R_{ij}|}.
\end{equation}

The sum covers every grid cell, including zeros outside the WUI update band;
nodata and malformed rasters are rejected rather than omitted.  HRRPUA and
both heat fields use the same metric over all dump times and cells, which also
detects an unintended secondary source and requires the final transient maps
to be zero after ELMFIRE's step-level reset.  A 0.5-percent map tolerance permits single-
precision raster output while remaining much smaller than any intended branch
or parameter response.  Paired algebraic invariants use $10^{-5}$; rotation
uses the primary 0.5-percent tolerance because trigonometric evaluation occurs
in single precision. Each simulation condition must have a confirmed successful
completion with matching SHA-256 identities for the simulation settings,
spatial inputs, both fuel-property tables, reference and analysis methods,
ELMFIRE program, simulation record, and all results. Because ELMFIRE's first
scheduled dump occurs after one scheduled output interval, the dump simulation specification
must contain exactly 60 and
120 seconds, with a unique final record at the
configured stop. Output labels alone cannot establish a valid simulation.
The model-program identity is retained, but its relationship to the stated
source revision must be established separately.

\begin{center}

\begin{longtable}{@{}P{0.23\textwidth}P{0.30\textwidth} P{0.18\textwidth}P{0.15\textwidth}@{}}
\toprule
Expected property & Observable and metric & Criterion & Failure sensitivity\\
\midrule
Piecewise source curve & All HRR transient rasters; normalized $L_1$ & $\le0.005$ & timing or source contamination\\
Ellipse heat kernel & Full DFC and radiation maps at peak HRRPUA & each $\le0.005$ & branches, coefficients, orientation\\
$D$ clamp and input paths & Pairwise full-map normalized $L_1$ & each $\le10^{-5}$ & clamp or I/O divergence\\
Rotation covariance & Rotated DFC-plus-radiation maps & $\le0.005$ & angle, sign, or raster orientation\\
Vegetation-source factor & Full DFC/radiation maps for the forest-source simulation condition & each $\le0.005$ & unused 1/3 factor or initial-source HRR suppression\\
No propagation & Final time of arrival (TOA) and off-source HRR cells & zero additional cells & confounding source creation\\
\bottomrule
\end{longtable}
\end{center}

All 23 simulation conditions and all required fields must be present. Overall PASS requires
every metric to pass; missing, duplicate, malformed, stale, or mistimed
evidence produces NOT EVALUABLE.

\section{Simulation configuration and predicted outcome}

\begin{center}

\begin{longtable}{@{}P{0.22\textwidth}P{0.30\textwidth}P{0.42\textwidth}@{}}
\toprule
Parameter & Configured value & Isolation role\\
\midrule
Grid & \CaseGridSize, \CaseCellSize\ m cells & full heat-kernel domain\\
CRS & EPSG:32610 & explicit metric coordinates\\
Time & $\Delta t=1$ s, stop 120 s, dump 60 s & exact peak-HRR dump\\
Initial source & one central fire behavior fuel model (FBFM) 91 cell & stationary WU-E source\\
Background & nonburnable FBFM93, adjustment 0 & prevents propagation\\
Spotting/terrain & disabled/flat & removes confounders\\
Model selector & WUI model type 2 & UMD/UCB implementation\\
Ellipse adjustment & 1.0 & intended default effective-flame factor\\
\bottomrule
\end{longtable}
\end{center}

The generated constant fields and rasters use identical north-up georeferencing,
nodata, dimensions, and data types.  The point ignition scheduled at 121 s is
deliberately after the 120 s stop; it prevents generic stalled-front termination
without adding a physical source during the experiment.

\begin{figure}[p]
  \centering
  \includegraphics[width=\textwidth]{input_configuration.pdf}
  \caption{Whole-domain prepared configuration generated from the
  the 270-degree wind-direction condition GeoTIFF inputs.  The single negative PHI cell is the
  initial WU-E source; FBFM93 and zero adjustment prevent secondary spread.}
\end{figure}

The predicted output is a fixed-source HRR history and wind-oriented DFC and
radiation fields.  Branch-adjacent simulation conditions may jump because the source itself
uses distinct regressions; the test checks that behavior rather than imposing
continuity absent from the implementation.

\section{Actual simulation results}

The current analysis procedure state is
\textbf{\WorkflowStatus}; the scientific decision is
\textbf{\VerificationStatus}, with \CompletedVariantCount{} of
\RequiredVariantCount{} required simulation conditions having evaluable execution evidence.
Reason: \VerificationReason.  The common executed-binary SHA-256 is
\ReadableHash{\ExecutedBinarySHA}.

\IfFileExists{../figures/heat_map_comparison.pdf}{%
\begin{figure}[p]
  \centering
  \includegraphics[width=\textwidth]{heat_map_comparison.pdf}
  \caption{Representative whole-domain ELMFIRE result, independent reference calculation, and
  difference.  This figure is generated only from completed case-specific output.}
\end{figure}
}{\noindent\fbox{\parbox{0.96\textwidth}{Actual-result figure unavailable because the
required ELMFIRE simulation conditions have not been evaluated.}}}

\section{Calculated metrics and verification decision}

\begin{longtable}{@{}P{0.19\textwidth}P{0.76\textwidth}@{}}
\toprule
\multicolumn{2}{@{}P{0.97\textwidth}@{}}{\textbf{Maximum HRRPUA normalized L1 error} --- \textbf{NOT EVALUABLE}}\\*
Criterion & \ensuremath{\leq} 0.005 \\*
Observed & Not available \\[0.6em]
\midrule
\multicolumn{2}{@{}P{0.97\textwidth}@{}}{\textbf{Maximum DFC full-map normalized L1 error} --- \textbf{NOT EVALUABLE}}\\*
Criterion & \ensuremath{\leq} 0.005 \\*
Observed & Not available \\[0.6em]
\midrule
\multicolumn{2}{@{}P{0.97\textwidth}@{}}{\textbf{Maximum radiation full-map normalized L1 error} --- \textbf{NOT EVALUABLE}}\\*
Criterion & \ensuremath{\leq} 0.005 \\*
Observed & Not available \\[0.6em]
\midrule
\multicolumn{2}{@{}P{0.97\textwidth}@{}}{\textbf{D greater than 50 m clamp pair error} --- \textbf{NOT EVALUABLE}}\\*
Criterion & \ensuremath{\leq} 1e-05 \\*
Observed & Not available \\[0.6em]
\midrule
\multicolumn{2}{@{}P{0.97\textwidth}@{}}{\textbf{Constant/raster parameter equivalence error} --- \textbf{NOT EVALUABLE}}\\*
Criterion & \ensuremath{\leq} 1e-05 \\*
Observed & Not available \\[0.6em]
\midrule
\multicolumn{2}{@{}P{0.97\textwidth}@{}}{\textbf{Maximum 90-degree covariance error} --- \textbf{NOT EVALUABLE}}\\*
Criterion & \ensuremath{\leq} 0.005 \\*
Observed & Not available \\[0.6em]
\midrule
\multicolumn{2}{@{}P{0.97\textwidth}@{}}{\textbf{Vegetation-source fuel-factor heat error} --- \textbf{NOT EVALUABLE}}\\*
Criterion & \ensuremath{\leq} 0.005 \\*
Observed & Not available \\[0.6em]
\midrule
\multicolumn{2}{@{}P{0.97\textwidth}@{}}{\textbf{Maximum arrival-mask mismatch count} --- \textbf{NOT EVALUABLE}}\\*
Criterion & 0 \\*
Observed & Not available \\[0.6em]
\midrule
\multicolumn{2}{@{}P{0.97\textwidth}@{}}{\textbf{Maximum unexpected HRR source count} --- \textbf{NOT EVALUABLE}}\\*
Criterion & 0 \\*
Observed & Not available \\[0.6em]
\midrule
\bottomrule
\end{longtable}


\textbf{Current decision: \VerificationStatus.}  NOT RUN is a analysis procedure state,
not evidence of verification and never maps to PASS.

\section{Technical interpretation}

The current result is limited by missing time-history evidence, rather than a
calculated heat-map discrepancy. The solver's low-tagged-node stopping rule
terminates this isolated-source setup before the required scheduled heat dumps.
The post-stop ignition disables a different, stalled-PHI shortcut; it does not
bypass the low-node rule. The lone overrun dump cannot replace the missing
history and is not interpreted as a measurement at the configured stop time.

The intended heat kernel also differs from the reviewed source in its angular
sectors, heat fractions, wind efficiency, and spatial normalization. These
source differences identify what a complete future run should test; they are
not numerical verification results from the incomplete run. After the source
can sustain a stationary thermal calculation, passing HRR with failing heat
maps would isolate the discrepancy to ellipse geometry or heat transfer.

\section{Scope and limitations}

This case covers the deterministic model-2 source ellipse and transient heat
distribution for a stationary structure.  It does not verify physical
calibration, propagation speed, ignition thresholds, embers, interface
transition, fuel consumption, or Statewide WU-E.  It intentionally treats
the pinned source as the implementation under test and excludes the distinct
Statewide WU--E formulation. CASE44 addresses the
HRR and heat-parameter response not isolated here.

\begin{thebibliography}{9}
\bibitem{PurnomoPCI2024} D. M. J. Purnomo et al., ``Reconstructing modes of destruction in wildland--urban interface fires using a semi-physical level-set model,'' \emph{Proceedings of the Combustion Institute}, 40 (2024), 105755. \href{https://doi.org/10.1016/j.proci.2024.105755}{doi:10.1016/j.proci.2024.105755}.
\bibitem{PurnomoIJWF2024} D. M. J. Purnomo et al., ``Integrating an urban fire model into an operational wildland fire model to simulate one dimensional wildland--urban interface fires: a parametric study,'' \emph{International Journal of Wildland Fire}, 33 (2024), WF24102. \href{https://doi.org/10.1071/WF24102}{doi:10.1071/WF24102}.
\end{thebibliography}
